% Prologue of Deep Learning as an Experimental Science.
\specialchapter{Prologue}{Ten runs of the same algorithm}{0}\label{ch:prologue}
\begin{refsection}[chapters/00-prologue.bib]
\labfloatprefix{P}

In 2017 Peter Henderson and five colleagues trained one reinforcement-learning algorithm, TRPO, on one simulated task, HalfCheetah, ten times. Nothing changed between the ten runs except the random seed. They split the runs into two groups of five and averaged the learning curves of each group. The two average curves came out visibly apart, and a two-sample $t$-test computed across the training curves gave $t=-9.09$ and $p=0.0016$ \parencite{HendersonEtAl2018}. Had the two groups been labelled ``baseline'' and ``our method'', the figure would have looked like a result.

Nothing in that experiment was broken. The algorithm was the same in both groups, the code was the same, and the arithmetic of the test was correct. What made the comparison misleading was its design: five runs per arm of a method whose runs vary a great deal, compared by a test whose notion of an observation did not match the experiment. Henderson's group ran it deliberately, as a control, at a time when published comparisons in deep reinforcement learning commonly rested on fewer than five runs per method. It is a small version of every chapter of this book.

\addsec{The same experiment on a laptop}

The phenomenon needs neither reinforcement learning nor a cluster. The scripts of this book train small networks on the handwritten digits that ship with scikit-learn, 1797 images of eight by eight pixels, in a fraction of a second per run.\sidenote{Every experiment in the book is a Python script in \texttt{figures/scripts}, built on a shared kit, \script{labkit.py}, and runs on a laptop in minutes. The prologue's is \script{00-prologue-split.py}.} Ten runs of one configuration, with a fifth of the training labels replaced by random ones to make the runs noisier, reach final test accuracies between $0.937$ and $0.961$, with a standard deviation of $0.007$ across runs. There are $126$ ways to split ten runs into two groups of five. For each split, compare the two groups in two ways.

The first way treats each run as one observation: summarize every run by its average test accuracy over the second half of training, and compare the five numbers of one group with the five of the other by Welch's $t$-test. Under this test, $2.4$ percent of the splits give $p<0.05$, about what a valid test of two identical methods should give, and none gives $p<0.001$.

The second way pools the points of the learning curves, $30$ evaluations per run in the second half of training, and compares the $150$ points of one group with the $150$ of the other as if each were an independent observation. Under this test, $54$ percent of the splits give $p<0.05$ and $29$ percent give $p<0.001$; the most extreme split gives $p=8\times10^{-11}$, while the run-level test on the same split gives $p=0.016$ (\cref{fig:prologue:split}). Evaluations of one run are not independent observations of the method: they share that run's initialization, data order and luck, and pooling them multiplies the apparent sample size by thirty without adding information about how the method varies from run to run.

\begin{figure}[t]
\begin{fullwidth}
\centering
\begin{tikzpicture}
\begin{groupplot}[group style={group size=2 by 1,horizontal sep=19mm},lplot,width=0.47\linewidth,height=0.36\linewidth]
\nextgroupplot[xlabel={training step},ylabel={test accuracy},ymin=0.5,ymax=1,xmin=0,xmax=1180,minor tick num=1,legend pos=south east]
\foreach \r in {0,...,4}{\addplot[inkblue!40,line width=0.4pt,no marks,forget plot] table[x=step,y=a\r] {figures/data/00-prologue-curves.dat};}
\foreach \r in {0,...,4}{\addplot[rubric!40,line width=0.4pt,no marks,forget plot] table[x=step,y=b\r] {figures/data/00-prologue-curves.dat};}
\addplot[inkblue,line width=1.2pt,no marks] table[x=step,y=meana] {figures/data/00-prologue-curves.dat};
\addplot[rubric,line width=1.2pt,no marks] table[x=step,y=meanb] {figures/data/00-prologue-curves.dat};
\legend{group 1 (5 runs),group 2 (5 runs)}
\nextgroupplot[xmode=log,xlabel={$p$-value},ylabel={share of the 126 splits},xmin=1e-11,xmax=1,ymin=0,ymax=1,grid=major,xtick={1e-10,1e-8,1e-6,1e-4,1e-2,1},legend pos=north west]
\addplot[reference,domain=1e-11:1,samples=80,forget plot] {x};
\addplot[reference,forget plot] coordinates {(0.05,0) (0.05,1)};
\addplot[armA,mark=none,const plot] table[x=prun,y=rank] {figures/data/00-prologue-pvalues.dat};
\addplot[armB,mark=none,const plot] table[x=ppool,y=rank] {figures/data/00-prologue-pvalues.dat};
\legend{one outcome per run,curve points pooled}
\end{groupplot}
\end{tikzpicture}
\end{fullwidth}
\caption{Ten runs of one configuration of a small network on the scikit-learn digits, with 20\% of the training labels randomized, differing only in their seeds. Left: the split into two groups of five that looks most different when curve points are pooled; thin lines are runs, thick lines group means. Right: for all $126$ splits, the share whose $p$-value is at most the value on the horizontal axis, for Welch's test on one outcome per run (ink) and for a $t$-test that pools the $30$ late evaluation points of each run as independent observations (vermilion). The dashed diagonal is what a valid test of two identical methods gives; the vertical line marks $p=0.05$. Script \script{00-prologue-split.py}.}\label{fig:prologue:split}
\end{figure}

Ecologists named this error more than forty years ago. In 1984 Stuart Hurlbert reviewed 176 field experiments published since 1960 and found that about a quarter of them, and nearly half of those that used inferential statistics, tested treatments against measurements that were not independent replicates: several samples from one plot, one tank or one animal, or repeated measurements of one unit over time. He called the practice \term{pseudoreplication}\index{pseudoreplication}, and the second form, which is what pooling the points of a learning curve amounts to, temporal pseudoreplication \parencite{Hurlbert1984}. The cure is to identify the \term{experimental unit}\index{experimental unit}, the smallest thing to which the treatment is independently applied, and to count replicates in those units. In a training experiment the unit is almost always the run.

\addsec{An experimental science without a laboratory course}

Most of what the field of deep learning knows, it knows by experiment. That one optimizer works better than another, that a given compute budget should be split between model size and data in a given proportion, that a normalization layer or a learning-rate schedule helps, that large models can do things small ones cannot: these are experimental findings, and theory, when it comes, usually comes after. In this sense deep learning is an experimental science, and a young one. Its experiments are expensive, noisy and easy to vary, its results travel fast, and its practitioners are trained mostly in mathematics, programming and engineering.

Older experimental sciences teach method before they teach results. A physics student spends a semester in a teaching laboratory measuring things whose values are known, learning to attach an uncertainty to every number and to propagate it \parencite{Taylor1997}. A biology student learns to include controls, to blind the observer, and to replicate. The design of experiments, as a branch of statistics, grew up around the problem of drawing conclusions from noisy and expensive runs. Ronald Fisher joined the Rothamsted agricultural station in 1919, where a single run of an experiment was a growing season, and the methods he developed there, replication, randomization and blocking, became \emph{The Design of Experiments} \parencite{Fisher1935}. A training run of a large network is, in cost and in noise, closer to a growing season than to a laboratory titration, and the experimental designs that Fisher and his successors developed for fields and factories \parencite{BoxHunterHunter2005} apply to it almost unchanged.

The method these courses teach exists because the alternative has a history. In 1903 Ren\'e Blondlot, a respected physicist at Nancy, announced a new kind of radiation, the N-rays, and within a year dozens of papers had reported their properties. In September 1904 the American physicist Robert Wood visited Blondlot's laboratory and, in the dark, quietly removed the aluminum prism that was supposed to disperse the rays into a spectrum. Blondlot went on reporting the spectrum. Wood's letter to \emph{Nature} ended the matter for most physicists \parencite{Wood1904}. The observers had not lied; they had seen what they expected to see, and the experiment had no control that could tell them otherwise. Blinding exists because of episodes like this one.

Deep learning has no laboratory course. Its experimental method is passed on by word of mouth, in reading groups and code reviews, and unevenly. How much seed variance is enough to matter? How can an unequal tuning budget flip a comparison? How should a scaling law be fitted, and how far can it be trusted? How should an ablation be designed so that it does not confound? The field has answered each of these questions somewhere, usually in a paper that showed a famous conclusion to be wrong, and has written down general warnings about its own practices \parencite{SculleyEtAl2018,LiptonSteinhardt2018}. The answers have not been collected into a course. Paul Cohen made the same complaint about artificial intelligence as a whole in 1995, and wrote a book on empirical methods for it \parencite{Cohen1995} before deep learning existed. A recent book and a widely used guide stand nearest to this one. Moritz Hardt's \emph{The Emerging Science of Machine Learning Benchmarks} treats the science of benchmarks and of the rankings they produce \parencite{Hardt2026}, and the \emph{Deep Learning Tuning Playbook} of Google researchers gives practical advice on tuning \parencite{GodboleEtAl2023}. This book is about the training experiment: how to run one, how to compare two, and how to know what a result means.

\addsec{The anatomy of a training experiment}

A claim about deep learning is usually a claim about what happens when something is changed: an optimizer, an architecture, a schedule, a dataset, a model size. The thing changed is the \term{treatment}\index{treatment}, and its versions are the \term{levels} of a \term{factor}\index{factor}. The \term{run}\index{run}, one training of one network from initialization to the end, is the experimental unit, and its \term{outcome}\index{outcome} is a number: a test error, a loss, an accuracy. Everything else that can change the outcome falls into one of two kinds. \term{Held-fixed variables}\index{held-fixed variable} are set to one value for every run: the dataset, the number of steps, the batch size, the hardware. \term{Nuisance variables}\index{nuisance variable} vary from run to run without being of interest: the seed of the initialization, the order of the data, the dropout masks, the nondeterminism of parallel hardware. A \term{budget}\index{budget} limits how many runs, how many hyperparameter trials, or how much compute the experiment may use, and a \term{protocol}\index{protocol} says how the runs are chosen, how their outcomes are summarized, and how the conclusion is drawn.

Each chapter records the design of the experiment it examines in a box like the following, which records Henderson's control experiment.

\begin{experiment}[Henderson et al.\ 2018, a control]
\field{question}{Do two groups of runs of the same algorithm differ?}
\field{treatment}{none: the group label was assigned by splitting the runs}
\field{held fixed}{algorithm (TRPO), task (HalfCheetah), hyperparameters, code}
\field{varied}{the random seed of each run}
\field{response}{the learning curve: return as a function of training steps}
\field{replicates}{five runs per group}
\field{analysis}{group-mean curves; a two-sample $t$-test across the training curves}
\end{experiment}

Two consequences of this vocabulary run through the book. The first is that a method is not a number. Run with a new seed, it produces a new outcome, and what a method has is a \term{run distribution}\index{run distribution}: the distribution of outcomes over everything that the experimenter does not control. A comparison of two methods is a comparison of two distributions, observed through a few samples. The second is that the variables held fixed are part of the claim. ``Method A beats method B'' always means ``at these settings of everything held fixed'', and a surprising number of famous conclusions turned out to be true only at the settings at which they were measured.

\addsec{Five questions}

The chapters are grouped by the question that each part teaches the reader to ask of an experiment.

\paragraph{How far does a number move when nothing changes?} Part I, \emph{Noise}, measures the run distribution: the spread that seeds produce and how many runs it takes to see past it (\cref{ch:seeds}), why fixing the seed does not fix a run (\cref{ch:instability}), and how to summarize a method over many tasks with few runs each without being fooled by the summary (\cref{ch:aggregates}).

\paragraph{Is the comparison fair?} Part II, \emph{Fair Comparison}, asks whether two arms differ only in what the claim is about: in their tuning budgets (\cref{ch:budgets}), in the presence of a placebo arm (\cref{ch:controls}), in the training recipe that came with a new architecture (\cref{ch:confounding}), and in what was held fixed while the treatment changed (\cref{ch:heldfixed}).

\paragraph{Is the cause identified?} Part III, \emph{Causes}, is about explanations: how to design an ablation (\cref{ch:ablations}), how to test a claim about mechanism by intervening on it (\cref{ch:mechanisms}), and how to read a correlation found across a sweep of runs (\cref{ch:sharpness}).

\paragraph{Does the result extend beyond what was measured?} Part IV, \emph{Scale}, treats scaling laws as measurements: how to fit one and state its uncertainty (\cref{ch:fitting}), how the protocol of a scaling study shapes the exponent it finds (\cref{ch:protocol}), and how far a trend can be extrapolated (\cref{ch:extrapolation}).

\paragraph{What does the number measure?} Part V, \emph{Measurement}, examines the instrument: how the choice of a metric can create an apparent phenomenon (\cref{ch:emergence}), and how the settings of an evaluation, such as the wording of a prompt, move its readings (\cref{ch:instruments}).

\addsec{How each chapter is built}

Every chapter opens with a famous conclusion, quoted as the field received it, with its source in the margin and a one-line verdict beneath it. The chapter then performs the conclusion's post-mortem. \emph{The claim} says who claimed what, on what evidence, and why it was reasonable to believe. \emph{The experiment} reconstructs the original design in a record like the one above. \emph{The post-mortem} reports what later work found and names the mechanism of the failure, and ends with a verdict, stamped in the margin, from a fixed vocabulary (\cref{tab:prologue:verdicts}). \emph{The method} teaches the general technique that would have caught the failure, with the statistics or mathematics that make it precise and at least one result proved completely. \emph{In the laboratory} reproduces the failure and its fix in a small experiment that runs on a laptop. \emph{Protocol} turns the method into a checklist for the next experiment of the same kind. Exercises follow, of five kinds: consolidation of the chapter's results, laboratory work with its scripts, the design of new experiments, the critique of described ones, and research problems.

\begin{table}[t]
\centering\small
\renewcommand{\arraystretch}{1.15}
\begin{tabularx}{\linewidth}{@{}>{\scshape\color{rubric}}l X@{}}
\toprule
reversed & under a fair protocol the comparison or conclusion flips\\
vanished & the effect is within noise, or disappears under a fair protocol\\
confounded & the effect is real but belongs to another variable that changed with the treatment\\
artifact & the effect is produced by the measurement: the metric, the fit or the evaluation\\
narrowed & the claim holds only in a smaller regime than claimed\\
misattributed & the effect is real; the explanation offered for it is not\\
unreplicated & an independent reanalysis or rerun does not reproduce the reported result\\
open & the evidence does not yet settle the question\\
\bottomrule
\end{tabularx}
\caption{The verdicts of the post-mortems.}\label{tab:prologue:verdicts}
\end{table}

The post-mortems are \emph{blameless}, in the sense that the term has in engineering, where the review of an outage looks for the conditions that made the failure possible rather than for the person who made it \parencite{BeyerEtAl2016}. Most of the conclusions examined here were reasonable given the evidence available when they were drawn, several were stated more carefully by their authors than by the field that repeated them, and several of the post-mortems were written by the original authors or their colleagues. A verdict applies to a claim, never to a person. The point of studying a famous wrong conclusion is that the experiment behind it was usually done by able people following the practice of their time, which is exactly the situation of the reader.

A laboratory notebook has one rule about mistakes: never erase. A wrong entry is struck through once, so that it stays legible, and the correction is written beside it. This book treats the field's famous conclusions in the same way. They are kept, stated fairly and legibly, and each is followed by what was learned from it.

\addsec{How to use the laboratory}

The laboratory sections are meant to be run. Every figure built from training runs is produced by a script in \texttt{figures/scripts}, which writes its data to \texttt{figures/data} and prints the numbers quoted in the text. The scripts need Python with NumPy, SciPy, scikit-learn and PyTorch, and each runs on a laptop's processor in at most a few minutes. They share a small kit, \script{labkit.py}, that trains a network on the scikit-learn digits or on a synthetic regression task and keeps four sources of randomness apart: the initial weights, the order of the data, the split into training and test sets, and the noise of labels, dropout and augmentation. Each lab sheet ends with a question to answer \emph{before} running the script. The prediction is the point: an experimenter who has written down a prediction is harder to fool, by the data or by their own hopes, than one who has not.\sidenote{A prediction written down in advance is the informal version of a preregistered analysis plan. The epilogue returns to it.}

The ornament at the foot of each chapter's opening page is a laboratory record too. It plots sixteen draws from a normal distribution, one more for each chapter: the open circles are the draws, the filled ones their running mean, and the bars a $95$ percent confidence interval for the mean, computed from the draws so far. The dashed line is the true mean. The first draw lands far above it, and there is no interval yet, because one run cannot estimate its own variance; with two or three draws the interval runs off the page. By the last chapter it is narrow and the mean sits on the line. It is the first lesson of \cref{ch:seeds}, drawn in the margin of the whole book.

\notesandrefs

\Textcite{HendersonEtAl2018} examined reproducibility in deep reinforcement learning more broadly: the effect of hyperparameters, network architecture, reward scale, random seeds, environments and codebases on reported results, and the statistical tests appropriate to them; their seed experiment is in the section on random seeds and trials. \Textcite{ColasSigaudOudeyer2018} turned its lesson into a tutorial on statistical power for the same setting, and \textcite{AgarwalEtAl2021} into a set of reporting practices; both return in Part I.

Hurlbert's paper on pseudoreplication \parencite{Hurlbert1984} is widely cited and readable without a background in ecology. Fisher's \emph{Design of Experiments} \parencite{Fisher1935} builds randomization, replication and the null hypothesis from the experiment of a lady who claims to tell by taste whether the milk or the tea was poured first into her cup; Box, Hunter and Hunter \parencite{BoxHunterHunter2005} is the standard modern text on designed experiments, and Taylor \parencite{Taylor1997} the standard introduction to error analysis in a physics laboratory. Wood's letter on the N-rays is short and worth reading in the original \parencite{Wood1904}.

The practice of blameless post-mortems is described in the chapter on postmortem culture of Google's book on site reliability engineering \parencite{BeyerEtAl2016}. Two essays that a reader of this book should know are \textcite{SculleyEtAl2018}, on how the pace of the field strains its empirical rigor, and \textcite{LiptonSteinhardt2018}, on patterns in machine learning writing that blur what was shown.

\printchapterbib
\end{refsection}
\labfloatrestore
